\begin{finding}[\fid{F21}: the specified rank-three tensor lattice is realized]\label{f:r3tensor}
Use the ring $R_0$, algebra $A$, center $C$, roots $z_1,z_2$ and named
generators of \fid{F19}. Let $V=W=R_0^2$ and give
$M=(V\otimes W)^{\otimes3}$ the actual right tensor-monomial action of
GCT~IV, \S6.1, p.~33, and Definition~11.1, equation~(93). Write
$\epsilon_+$ for $P_i=D$, $\epsilon_-$ for $P_i=0$, and
\[
 N(a):\quad P_1=\begin{pmatrix}D&0\\a&0\end{pmatrix},\qquad
 P_2=\begin{pmatrix}0&a\\0&D\end{pmatrix}.
\]
These are right row matrices, as in IV Proposition~11.14,
equation~(110). There is an action-preserving isomorphism over $R_0$,
\begin{equation}\label{eq:r3tensor}
 M\simeq\epsilon_+^{20}\oplus\epsilon_-^4
       \oplus N(\delta)^{16}\oplus N(D-2)^4.
\end{equation}
In particular every displayed change of basis is regular at $v=1$ in
characteristic different from $2,3$; no colliding-root difference is inverted.

Put $L_i=C/(z-z_i)$ and
$S=\Spec R_0\sqcup\Spec R_0\sqcup\Spec C$. For
\[
 E=(\mathcal O,\mathcal O,\mathcal O_C^2),\qquad
 \mathcal F=(\mathcal O^{20},\mathcal O^4,
             \widetilde L_1^{16}\oplus\widetilde L_2^4),
\]
the coherent model realizes this same action as
\[
 M\simeq\Gamma(S,\mathcal Hom(E,\mathcal F)),\qquad A=\Gamma(S,\End E),
\]
with right action by precomposition. It is faithful before specialization
and kills the four-dimensional radical after $v=1$.
\end{finding}
\begin{proof}
Let $L=R_0^2$. The cited ordinary tensor action, with the consistent
normalization $c_i=T_i+v^{-1}$, is
\[
 \mathbf v_{\mathbf k}c_i=
 \begin{cases}
 v\mathbf v_{\mathbf k}+\mathbf v_{\mathbf k s_i},&k_i<k_{i+1},\\
 \delta\mathbf v_{\mathbf k},&k_i=k_{i+1},\\
 v^{-1}\mathbf v_{\mathbf k}+\mathbf v_{\mathbf k s_i},&k_i>k_{i+1}.
 \end{cases}
\]
On $(112,121,211)$ it gives
\[
 c_1=\begin{pmatrix}\delta&0&0\\0&v&1\\0&1&v^{-1}\end{pmatrix},\qquad
 c_2=\begin{pmatrix}v&1&0\\1&v^{-1}&0\\0&0&\delta\end{pmatrix}.
\]
The row-basis matrix
\[
 B_3=\begin{pmatrix}v^2&v&1\\-1-v^{-2}&v&1\\
                         -v^{-1}&-v^{-2}&\delta\end{pmatrix}
\]
has determinant $(D-1)^2$. Multiplying gives
\[
 B_3c_1B_3^{-1}=\begin{pmatrix}\delta&0&0\\0&\delta&0\\0&1&0\end{pmatrix},
 \quad
 B_3c_2B_3^{-1}=\begin{pmatrix}\delta&0&0\\0&0&1\\0&0&\delta\end{pmatrix}.
\]
Thus this block is $\epsilon\oplus U$, with
\[
 c_1|_U=\begin{pmatrix}\delta&0\\1&0\end{pmatrix},\qquad
 c_2|_U=\begin{pmatrix}0&1\\0&\delta\end{pmatrix}.
\]
On $(122,212,221)$ the original $c_1,c_2$ matrices interchange; applying
$B_3$ and swapping its last two rows again gives $\epsilon\oplus U$.
The weight spaces $111$ and $222$ are each $\epsilon$. Since the determinants
are units, this proves $L^{\otimes3}=\epsilon^4\oplus U^2$ over $R_0$,
not just over its fraction field.

By $P_i=c_i\otimes c_i+(c_i-\delta)\otimes(c_i-\delta)$,
$\epsilon\otimes\epsilon=\epsilon_+$, and each mixed tensor is
$N(\delta)$. On $U\otimes U$ in the order $(11,12,21,22)$,
\[
 P_1=\begin{pmatrix}D&0&0&0\\\delta&0&0&0\\\delta&0&0&0\\
                         2&-\delta&-\delta&D\end{pmatrix},\quad
 P_2=\begin{pmatrix}D&-\delta&-\delta&2\\0&0&0&\delta\\
                         0&0&0&\delta\\0&0&0&D\end{pmatrix}.
\]
Take
\[
 B_4=\begin{pmatrix}
 \delta&-1&-1&\delta\\0&-1&1&0\\
 2-D&-\delta&-\delta&D\\-D&\delta&\delta&D-2
 \end{pmatrix},\qquad \det B_4=8(D-1)^2.
\]
Direct multiplication yields
\[
 B_4P_1B_4^{-1}=\begin{pmatrix}D&0&0&0\\0&0&0&0\\0&0&D&0\\0&0&D-2&0\end{pmatrix},
\quad
 B_4P_2B_4^{-1}=\begin{pmatrix}D&0&0&0\\0&0&0&0\\0&0&0&D-2\\0&0&0&D\end{pmatrix}.
\]
This is $\epsilon_+\oplus\epsilon_-\oplus N(D-2)$, again over $R_0$.
In $(\epsilon^4\oplus U^2)\otimes(\epsilon^4\oplus U^2)$ the unmixed
trivial contribution has multiplicity $16$, the two mixed contributions
have total multiplicity $16$, and $U\otimes U$ occurs four times.
This proves \eqref{eq:r3tensor}, whose rank is $20+4+32+8=64$.

It remains to match the named action to the center model, rather than
merely match dimensions. For $a=\delta,D-2$, put $q_a=a/D$, a unit, and
\[
 S_a=\begin{pmatrix}1&0\\-q_a^2&q_a\end{pmatrix}.
\]
The matrices of $S_a(P_1/D)S_a^{-1}$ and $S_a(P_2/D)S_a^{-1}$ are
respectively
\[
 \begin{pmatrix}1&0\\0&0\end{pmatrix},\qquad
 \begin{pmatrix}1-z_i&1\\z_i(1-z_i)&z_i\end{pmatrix},
 \quad z_i=1-q_a^2.
\]
They are exactly the standard row module for $M_2(C)$ evaluated on branch
$i$. The right-module formula with $\mathcal Hom$ follows. The two branch
evaluations jointly embed $C$ into $R_0^2$, so their modules have joint
annihilator zero; both scalar blocks also occur. At $v=1$ the branch
modules become the same residue module of $k[\varepsilon]/\varepsilon^2$,
and the central nilpotent acts by zero. Its matrix ideal is the radical
of \fid{F19}, proving the last assertion.
\end{proof}
The generic multiplicities themselves are already in IV \S12.3; the
additional assertion here is the displayed regular action isomorphism.
The coherent sheaf $\mathcal F$ is not locally free on the collision
component: its two generic branch ranks are $16$ and $4$. This result
does not identify a simultaneous canonical left quantum-group action
or VIII's based tensor lattice with the monomial lattice.

